\begin{afterword}
\summaryhead

The post draws a conclusion from the posts before it: if you care only about controversial science, you will fill your head with garbage. The media report only the frontier, where findings are often based on the thinnest evidence and wrong half the time; settled science is never news. Controversies are problems hard enough to fool experts, and for outsiders they offer only the fun of picking a side, which a football game also offers.

A good textbook gives careful explanations, derivations, experiments, exercises and a reasonably good guarantee of truth. A press release gives a fake explanation of a result its writer did not understand, which probably will not replicate. Science is built on discoveries going back to Archimedes, so one should start at the beginning and not be embarrassed to read elementary textbooks. Jumping to the frontier is like trying to jump over the lower half of Everest. Controversies can matter, but science need not be controversial, or recent, to be interesting.

\responsehead

The post gives advice, and the advice depends on two claims about science in the news. Both are hedged, and both have support from later research.

The first is that newsworthy science is ``often'' wrong half the time. Nine years after the post, Dumas-Mallet and colleagues traced biomedical association studies that newspapers had reported, and found that fewer than half, 48.7 per cent, were confirmed by later meta-analyses. That is one field and one kind of study, but it is close to the post's figure.

The second is that press releases usually convey only a ``delusion of understanding'' of results that ``probably'' will not replicate. Sumner and colleagues, studying press releases issued by British universities in 2011, found each of three kinds of exaggeration in a third or more of them. That supports the charge of distortion. The post's explanation, that the writer of the release ``didn't understand'' the result, is another matter: the study measured what the releases said, not what their writers understood, and the post gives no evidence for it.

The rest of the post is attitude and should be read as such. The advice to read textbooks is hedged where it needs to be (``a reasonably good guarantee'' of truth), grants that some controversies matter, and makes its case with jokes: a football game, a jump over half of Everest, a brain turned to liquid. Archimedes is a fair example: \textsc{On Floating Bodies} holds the first known statement of the principle now named after Archimedes.

In short: sound advice to learn settled science before the frontier, based on hedged claims about science news that later studies of newspapers and press releases broadly support, with an explanation of the press releases' faults, that their writers did not understand the results, that the post does not support.
\end{afterword}
